\documentclass[11pt]{article}

\usepackage{amsmath,amssymb}
\usepackage{xurl}
\usepackage[hidelinks]{hyperref}
\setlength{\emergencystretch}{2em}

% Shared commands for notes in docs/paper-gaps/.

\DeclareUnicodeCharacter{2080}{\ifmmode{}_{0}\else\textsubscript{0}\fi}
\DeclareUnicodeCharacter{2081}{\ifmmode{}_{1}\else\textsubscript{1}\fi}
\DeclareUnicodeCharacter{2082}{\ifmmode{}_{2}\else\textsubscript{2}\fi}
\DeclareUnicodeCharacter{2099}{\ifmmode{}_{n}\else\textsubscript{n}\fi}

\newcommand{\Lean}{\textsc{Lean}}
\ExplSyntaxOn
\NewDocumentCommand{\leanid}{m}{
  \ifmmode
    \textsf{\detokenize{#1}}
  \else
    \group_begin:
    \urlstyle{sf}
    \tl_set:Nn \l_tmpa_tl {#1}
    \tl_replace_all:Nnn \l_tmpa_tl {\_} {_}
    \exp_args:NV \path \l_tmpa_tl
    \group_end:
  \fi
}
\ExplSyntaxOff
\providecommand{\tr}{\operatorname{tr}}
\newcommand{\tnleanIssueBase}{https://github.com/LionSR/TNLean/issues/}
\newcommand{\tnleanPrBase}{https://github.com/LionSR/TNLean/pull/}
\newcommand{\tnleanDocBase}{https://sirui-lu.com/TNLean/docs/}
\newcommand{\ghissue}[1]{\href{\tnleanIssueBase#1}{GitHub issue \##1}}
\newcommand{\ghpr}[1]{\href{\tnleanPrBase#1}{PR \##1}}
\newcommand{\doclink}[1]{\href{\tnleanDocBase#1.html}{\path{#1}}}

% Dirac notation and the identity operator, as in blueprint/src/macros/common.tex.
% \providecommand keeps note-local definitions authoritative.
\usepackage{dsfont}
\providecommand{\ket}[1]{|#1\rangle}
\providecommand{\bra}[1]{\langle #1|}
\providecommand{\braket}[2]{\langle #1 | #2 \rangle}
\providecommand{\ketbra}[2]{|#1\rangle\!\langle #2|}
\providecommand{\Id}{\mathds{1}}
\providecommand{\one}{\mathds{1}}
\providecommand{\C}{\mathbb{C}}
\providecommand{\MN}[1]{M_{#1}(\C)}
\providecommand{\MD}{M_D(\C)}

% External referencing for paper sources.  The build helper writes sanitized
% auxiliary files under build/paper-aux/ and excludes duplicated source labels.
\usepackage{xr}

% Import a generated paper auxiliary file with a caller-supplied label prefix.
% The candidate paths support builds from docs/paper-gaps/, the repository root,
% and build/paper-aux/ itself.
\newcommand{\importpaperaux}[2]{%
  \IfFileExists{../../build/paper-aux/#2.aux}{%
    \externaldocument[#1]{../../build/paper-aux/#2}%
  }{%
    \IfFileExists{build/paper-aux/#2.aux}{%
      \externaldocument[#1]{build/paper-aux/#2}%
    }{%
      \IfFileExists{#2.aux}{%
        \externaldocument[#1]{#2}%
      }{}%
    }%
  }%
}

% General external reference macros
\newcommand{\paperref}[2]{\ref{#1-#2}}
\newcommand{\papereqref}[2]{(\ref{#1-#2})}

% CPSV16 source (1606.00608 / MPDO-22-12-17-2.tex) external document and shortcuts
\importpaperaux{cpsv16-}{1606.00608/MPDO-22-12-17-2-tnlean}
\newcommand{\cpsvref}[1]{\paperref{cpsv16}{#1}}
\newcommand{\cpsveqref}[1]{\papereqref{cpsv16}{#1}}

% Verdict marker: \gapnote{<kind>}{<status>}. Kind is the policy
% classification (clarification, local-correction, scope-restriction,
% unfaithful, false-source, open-gap); status is open, wip (elimination
% underway), resolved, or historical. Machine-read by the site index;
% prints a verdict line.
\newcommand{\gapnote}[2]{\par\noindent\textbf{Verdict:} #1 (#2).\par}


\title{Closure Record for the Strong-Subadditivity Equality
Characterization and Hayashi--Markov Decomposition}
\author{}
\date{2026-07-28}

\begin{document}
\maketitle
\gapnote{open-gap}{resolved}

\section{The assertion and its current formal status}

Let $\rho_{ABC}$ be a density operator on the finite-dimensional Hilbert space
$\mathcal H_A\otimes\mathcal H_B\otimes\mathcal H_C$, and let
$\rho_{AB}$, $\rho_{BC}$, and $\rho_B$ denote its reduced density operators.
Equality in strong subadditivity is the identity
\begin{align}
    S(\rho_{ABC})+S(\rho_B)
        &=S(\rho_{AB})+S(\rho_{BC}).
    \label{eq:cpsv16_ssa_equality_hayashi_markov_ssa_equality}
\end{align}
The standard structural characterization states that
Eq.~(\ref{eq:cpsv16_ssa_equality_hayashi_markov_ssa_equality}) holds if and
only if a unitary change of basis on $\mathcal H_B$ gives a finite orthogonal
direct-sum decomposition
\begin{align}
    \mathcal H_B
        &\cong\bigoplus_j
        \mathcal H_{B_j^L}\otimes\mathcal H_{B_j^R},
    \label{eq:cpsv16_ssa_equality_hayashi_markov_middle_decomposition}
\end{align}
under which
\begin{align}
    \rho_{ABC}
        &\cong\bigoplus_j
        p_j\,\rho_{AB_j^L}\otimes\rho_{B_j^RC}.
    \label{eq:cpsv16_ssa_equality_hayashi_markov_qmc_decomposition}
\end{align}
Here $p_j\geq 0$, $\sum_jp_j=1$, and $\rho_{AB_j^L}$ and
$\rho_{B_j^RC}$ are density operators on the indicated tensor factors.
Equation~(\ref{eq:cpsv16_ssa_equality_hayashi_markov_qmc_decomposition}) is
the quantum-Markov-chain structure of $\rho_{ABC}$; this characterization and
its relation to recovery maps are reviewed in Ref.~\cite{Cirac2021Matrix}.

Both directions of the equivalence between
Eqs.~(\ref{eq:cpsv16_ssa_equality_hayashi_markov_ssa_equality}) and
(\ref{eq:cpsv16_ssa_equality_hayashi_markov_qmc_decomposition}) are now
proved.\footnote{The forward implication is
\leanid{Matrix.hayashi_ssa_equality_characterization_forward} in
\path{TNLean/Analysis/EntropyMarkovForward.lean}, with the established root
name \leanid{hayashi_ssa_equality_characterization_forward} in
\doclink{TNLean/Entropy/SSAEqualityCharacterization}. The reverse implication
is \leanid{hayashi_ssa_equality_characterization_reverse} in
\path{TNLean/Analysis/EntropyMarkovReverse.lean}. The declaration
\leanid{hayashi_ssa_equality_characterization} combines the two directions.
The equality predicate is \leanid{IsSSAEquality} in
\doclink{TNLean/Analysis/Entropy}, and the structural predicate is
\leanid{HayashiMarkovDecomposition}.}
No axiom declaration occurs in this characterization.

The forward proof first derives the ambient Hayden--Jozsa--Petz--Winter block
identity
\begin{align}
    W^\dagger\rho_{ABC}W
        &=0_\perp\oplus\bigoplus_j\omega_j\otimes\widehat\rho_j,
    \label{eq:cpsv16_ssa_equality_hayashi_markov_ambient_hjpw}\\
    W
        &=\Id_A\otimes(U_B\otimes\Id_C).
    \label{eq:cpsv16_ssa_equality_hayashi_markov_ambient_unitary}
\end{align}
The matrices $\widehat\rho_j$ are the output states of the supported sectors.
They are positive, have trace one, and have the prescribed common-factor
marginals. Every complementary block is zero. The recovery construction orders
the middle factors as $B_j^R\otimes B_j^L$; reversing each fibre gives the
Hayashi order $B_j^L\otimes B_j^R$. The Hayashi unitary has the adjoint
orientation
\begin{align}
    V_B
        &=U_B^\dagger,
    \label{eq:cpsv16_ssa_equality_hayashi_markov_hayashi_unitary}\\
    V\rho_{ABC}V^\dagger
        &=W^\dagger\rho_{ABC}W.
    \label{eq:cpsv16_ssa_equality_hayashi_markov_unitary_orientation}
\end{align}

For every ambient sector, including each one-dimensional complementary sector,
define
\begin{align}
    p_j
        &=\operatorname{Re}\tr(\omega_j).
    \label{eq:cpsv16_ssa_equality_hayashi_markov_sector_weight}
\end{align}
Positivity and the ambient trace identity give
\begin{align}
    p_j
        &\geq 0,
    \label{eq:cpsv16_ssa_equality_hayashi_markov_weight_nonnegative}\\
    \sum_jp_j
        &=1.
    \label{eq:cpsv16_ssa_equality_hayashi_markov_weight_normalization}
\end{align}
On each supported sector, the right density operator is $\widehat\rho_j$. If
$p_j>0$, the left density operator is $\omega_j/p_j$. Total
positive-semidefinite normalization is applied only when the final
\leanid{HayashiMarkovDecomposition} is formed. If a supported sector has
$p_j=0$, this normalization selects a left filler while retaining
$\widehat\rho_j$ on the right. A complementary sector may use fillers on both
sides. In either case, the factor $p_j$ makes the fillers invisible. No filler
occurs in the raw identity
(\ref{eq:cpsv16_ssa_equality_hayashi_markov_ambient_hjpw}).

\section{The two proved implications}

\paragraph{Forward direction.}
Equality in strong subadditivity yields, in order, the product-reference
recovery identity on the supported input, the invariant conditional family,
its common invariant block form, the ambient bipartite decomposition, and the
pure-ancilla block action. Equations~(11), (14), and~(15) of the
Hayden--Jozsa--Petz--Winter argument, whose role in the equality
characterization is reviewed in Ref.~\cite{Cirac2021Matrix}, then give
Eq.~(\ref{eq:cpsv16_ssa_equality_hayashi_markov_ambient_hjpw}). The weight and
normalization argument in
Eqs.~(\ref{eq:cpsv16_ssa_equality_hayashi_markov_sector_weight})--%
(\ref{eq:cpsv16_ssa_equality_hayashi_markov_weight_normalization}) converts
this identity into
Eq.~(\ref{eq:cpsv16_ssa_equality_hayashi_markov_qmc_decomposition}). The formal
proof is \leanid{Matrix.hayashi_ssa_equality_characterization_forward}.

\paragraph{Reverse direction.}
For a block-diagonal state of the form
(\ref{eq:cpsv16_ssa_equality_hayashi_markov_qmc_decomposition}), the entropy
identities are
\begin{align}
    S\!\left(\bigoplus_jp_j\omega_j\right)
        &=H(\{p_j\})+\sum_jp_jS(\omega_j),
    \label{eq:cpsv16_ssa_equality_hayashi_markov_direct_sum_entropy}\\
    S(\omega\otimes\tau)
        &=S(\omega)+S(\tau).
    \label{eq:cpsv16_ssa_equality_hayashi_markov_tensor_entropy}
\end{align}
Applying
Eqs.~(\ref{eq:cpsv16_ssa_equality_hayashi_markov_direct_sum_entropy}) and
(\ref{eq:cpsv16_ssa_equality_hayashi_markov_tensor_entropy}) to the state and
its three marginals cancels the four entropy terms and proves
Eq.~(\ref{eq:cpsv16_ssa_equality_hayashi_markov_ssa_equality}). The formal
proof is \leanid{hayashi_ssa_equality_characterization_reverse}. The public
biconditional \leanid{hayashi_ssa_equality_characterization} depends only on
these two proved directions.

\section{Trace-preserving Petz channel}

Let $\sigma$ be a positive semidefinite reference matrix on a bipartite space,
and define its right-partial-trace marginal by
\begin{align}
    \tau
        &=\tr_R\sigma.
    \label{eq:cpsv16_ssa_equality_hayashi_markov_reference_marginal}
\end{align}
The Petz transpose map for the right partial trace is formalized on the support
of $\tau$ as
\begin{align}
    \mathcal R_\sigma(X)
        &=
        \sqrt{\sigma}
        \left(\tau^{-1/2}_{\mathrm{supp}}
            X
            \tau^{-1/2}_{\mathrm{supp}}
            \otimes\Id_R\right)
        \sqrt{\sigma}.
    \label{eq:cpsv16_ssa_equality_hayashi_markov_petz_support_map}
\end{align}
This map is completely positive in rectangular Kraus form, recovers the
reference matrix, and satisfies the support trace identity
\begin{align}
    \mathcal R_\sigma(\tau)
        &=\sigma,
    \label{eq:cpsv16_ssa_equality_hayashi_markov_petz_reference_recovery}\\
    \tr\!\left(\mathcal R_\sigma(X)\right)
        &=\tr(P_\tau X),
    \label{eq:cpsv16_ssa_equality_hayashi_markov_petz_support_trace}
\end{align}
where $P_\tau$ is the support projection of $\tau$.\footnote{Formal
declarations:
\leanid{Matrix.partialTraceRightPetzMap},
\leanid{Matrix.partialTraceRightPetzMap_isKrausCP},
\leanid{Matrix.partialTraceRightPetzMap_trace}, and
\leanid{Matrix.partialTraceRightPetzMap_partialTraceRight} in
\path{TNLean/Channel/PetzRecovery.lean}.}
Equation~(\ref{eq:cpsv16_ssa_equality_hayashi_markov_petz_support_map}) is the
support formula in Theorem~3, Equation~(8), of the
Hayden--Jozsa--Petz--Winter argument summarized in
Ref.~\cite{Cirac2021Matrix}.

For normalized $\sigma$, define
\begin{align}
    Q_\tau
        &=\Id-P_\tau,
    \label{eq:cpsv16_ssa_equality_hayashi_markov_support_complement}\\
    \omega_R
        &=\tr_L\sigma.
    \label{eq:cpsv16_ssa_equality_hayashi_markov_reference_right_marginal}
\end{align}
The complementary map is
\begin{align}
    \mathcal C_\sigma(X)
        &=Q_\tau XQ_\tau\otimes\omega_R.
    \label{eq:cpsv16_ssa_equality_hayashi_markov_petz_complement}
\end{align}
It is completely positive and contributes the trace
\begin{align}
    \tr\!\left(\mathcal C_\sigma(X)\right)
        &=\tr(Q_\tau X).
    \label{eq:cpsv16_ssa_equality_hayashi_markov_complement_trace}
\end{align}
Consequently,
\begin{align}
    \widehat{\mathcal R}_\sigma
        &=\mathcal R_\sigma+\mathcal C_\sigma
    \label{eq:cpsv16_ssa_equality_hayashi_markov_completed_petz_channel}
\end{align}
is completely positive and trace preserving. It agrees with
$\mathcal R_\sigma$ whenever $P_\tau XP_\tau=X$ and still satisfies
Eq.~(\ref{eq:cpsv16_ssa_equality_hayashi_markov_petz_reference_recovery}).%
\footnote{Formal declarations:
\leanid{Matrix.partialTraceRightPetzComplementMap},
\leanid{Matrix.partialTraceRightPetzChannel},
\leanid{Matrix.partialTraceRightPetzChannel_isKrausCPTP},
\leanid{Matrix.partialTraceRightPetzChannel_apply_of_supported}, and
\leanid{Matrix.partialTraceRightPetzChannel_partialTraceRight} in
\path{TNLean/Channel/PetzRecovery.lean}.}

The source support condition is connected to this completed channel. If
positive semidefinite matrices $\rho$ and $\sigma$ satisfy
$\ker\sigma\subseteq\ker\rho$, then their right marginals satisfy the same
kernel inclusion. With $\tau=\tr_R\sigma$, this gives
\begin{align}
    P_\tau(\tr_R\rho)P_\tau
        &=\tr_R\rho.
    \label{eq:cpsv16_ssa_equality_hayashi_markov_marginal_support}
\end{align}
The complementary term therefore vanishes on $\tr_R\rho$, so
\begin{align}
    \widehat{\mathcal R}_\sigma(\tr_R\rho)
        &=\mathcal R_\sigma(\tr_R\rho).
    \label{eq:cpsv16_ssa_equality_hayashi_markov_completed_raw_agreement}
\end{align}
\footnote{Formal declarations:
\leanid{Matrix.IsHermitian.supportProj_mul_mul_supportProj_of_mulVec_kernel_le}
and
\leanid{Matrix.partialTraceRightPetzChannel_apply_partialTraceRight_of_support}
in \path{TNLean/Channel/Schwarz/PetzRecoverySupport.lean}.}
Thus the choice of trace-preserving extension away from the support is no
longer part of the forward equality gap.

\section{Relative entropy and Weyl reduction}

Let $D(\rho\Vert\sigma)$ denote quantum relative entropy on its support domain.
If $\rho$ and $\sigma$ are positive semidefinite,
$\ker\sigma\subseteq\ker\rho$, and $\tau$ is positive semidefinite with
$\tr\tau=1$, then
\begin{align}
    D(\rho\otimes\tau\Vert\sigma\otimes\tau)
        &=D(\rho\Vert\sigma).
    \label{eq:cpsv16_ssa_equality_hayashi_markov_relative_entropy_ancilla}
\end{align}
The proof expands the logarithm of the singular tensor product through the
three support projections. Support absorption and $\tr\tau=1$ remove the
projection and ancilla terms.\footnote{Formal declaration:
\leanid{quantumRelativeEntropy_kronecker_support} in
\path{TNLean/Channel/Schwarz/PetzEqualityPrerequisites.lean}.}
This identity transports equality through the Weyl-twirl proof of data
processing, but it does not by itself show that equality implies recovery.

Fix a primitive $d_R$-th root of unity. For the corresponding Weyl operators
$W(a,b)$, define
\begin{align}
    U_{ab}
        &=\Id_L\otimes W(a,b),
    \label{eq:cpsv16_ssa_equality_hayashi_markov_weyl_unitary}\\
    \overline X
        &=\frac{1}{d_R^2}\sum_{c,e}U_{ce}XU_{ce}^\dagger,
    \label{eq:cpsv16_ssa_equality_hayashi_markov_weyl_average}\\
    \overline X
        &=(\tr_R X)\otimes(d_R^{-1}\Id_R).
    \label{eq:cpsv16_ssa_equality_hayashi_markov_weyl_partial_trace}
\end{align}
Suppose $\rho$ and $\sigma$ are positive semidefinite,
$\ker\sigma\subseteq\ker\rho$, and data processing under the right partial
trace is saturated. Then every $a,b$ satisfies
\begin{align}
    D(\overline\rho\Vert\overline\sigma)
        &=
        D(U_{ab}\rho U_{ab}^\dagger
            \Vert U_{ab}\sigma U_{ab}^\dagger).
    \label{eq:cpsv16_ssa_equality_hayashi_markov_weyl_summand_equality}
\end{align}
Equations~(\ref{eq:cpsv16_ssa_equality_hayashi_markov_weyl_partial_trace}) and
(\ref{eq:cpsv16_ssa_equality_hayashi_markov_relative_entropy_ancilla}),
together with the inherited support inclusion for the marginals, identify the
left-hand side with $D(\tr_R\rho\Vert\tr_R\sigma)$. Unitary invariance
identifies the right-hand side with $D(\rho\Vert\sigma)$.\footnote{Formal
declaration:
\leanid{quantumRelativeEntropy_weyl_average_eq_summand_of_partialTraceRight_eq}
in \path{TNLean/Channel/Schwarz/PetzEqualityPrerequisites.lean}.}
This result concerns scalar relative entropies only. It neither characterizes
equality in joint convexity nor supplies an operator intertwining or recovery
identity.

Summing
Eq.~(\ref{eq:cpsv16_ssa_equality_hayashi_markov_weyl_summand_equality}) with
uniform Weyl weights gives
\begin{align}
    D(\overline\rho\Vert\overline\sigma)
        &=
        \frac{1}{d_R^2}\sum_{c,e}
        D(U_{ce}\rho U_{ce}^\dagger
            \Vert U_{ce}\sigma U_{ce}^\dagger).
    \label{eq:cpsv16_ssa_equality_hayashi_markov_weyl_jensen_equality}
\end{align}
Unitary invariance makes every summand equal to $D(\rho\Vert\sigma)$, and the
$d_R^2$ weights sum to one.\footnote{Formal declaration:
\leanid{quantumRelativeEntropy_weyl_jensen_eq_of_partialTraceRight_eq} in
\path{TNLean/Channel/Schwarz/PetzEqualityPrerequisites.lean}.}
Again, this scalar equality does not characterize equality in joint convexity
and does not imply an operator intertwining or Petz recovery identity.

\section{Positive-definite resolvent argument}

For positive definite $\rho$ and $\sigma$, define the weighted Weyl family
\begin{align}
    A_g
        &=d_R^{-2}U_g\rho U_g^\dagger,
    \label{eq:cpsv16_ssa_equality_hayashi_markov_weighted_weyl_state}\\
    B_g
        &=d_R^{-2}U_g\sigma U_g^\dagger.
    \label{eq:cpsv16_ssa_equality_hayashi_markov_weighted_weyl_reference}
\end{align}
For $t>0$, define
\begin{align}
    T_g(X)
        &=A_gX+tXB_g,
    \label{eq:cpsv16_ssa_equality_hayashi_markov_local_left_right_operator}\\
    T(X)
        &=\left(\sum_gA_g\right)X
          +tX\left(\sum_gB_g\right),
    \label{eq:cpsv16_ssa_equality_hayashi_markov_summed_left_right_operator}\\
    \Lambda_t
        &=T^{-1}\!\left(\sum_gB_g\right).
    \label{eq:cpsv16_ssa_equality_hayashi_markov_common_resolvent_candidate}
\end{align}
Writing $R_g=B_g-T_g(\Lambda_t)$ and using the Hilbert--Schmidt inner product
gives the exact squared-defect identity
\begin{align}
    \sum_g
    \left\|
        T_g^{-1/2}(B_g)-T_g^{1/2}(\Lambda_t)
    \right\|_{\mathrm{HS}}^2
        &=
        \sum_g\operatorname{Re}
        \left\langle B_g,T_g^{-1}(B_g)\right\rangle_{\mathrm{HS}}
        \notag\\
        &\quad-
        \operatorname{Re}
        \left\langle
            \sum_gB_g,
            T^{-1}\!\left(\sum_gB_g\right)
        \right\rangle_{\mathrm{HS}}.
    \label{eq:cpsv16_ssa_equality_hayashi_markov_resolvent_defect_identity}
\end{align}
If the right-hand side vanishes, every squared norm on the left vanishes.
Therefore,
\begin{align}
    T_g^{-1}(B_g)
        &=\Lambda_t
    \label{eq:cpsv16_ssa_equality_hayashi_markov_common_resolvent}
\end{align}
for every $g$, including the identity Weyl element. This is the
Jen\v{c}ov\'a--Ruskai Appendix calculation in arXiv:0903.2895v4,
lines 1313--1343; the resulting common-resolvent equality occurs in
Section~4, lines 652--674. The relative-entropy equality methods underlying
this argument are reviewed in Ref.~\cite{Cirac2021Matrix}.\footnote{Formal
declarations:
\leanid{Matrix.weyl_resolvent_defect_identity} and
\leanid{Matrix.weyl_resolvent_eq_of_defect_eq_zero} in
\path{TNLean/Channel/Schwarz/PetzEqualityPrerequisites.lean}.}

Define
\begin{align}
    A
        &=\sum_gA_g,
    \label{eq:cpsv16_ssa_equality_hayashi_markov_summed_state}\\
    B
        &=\sum_gB_g.
    \label{eq:cpsv16_ssa_equality_hayashi_markov_summed_reference}
\end{align}
The positive-definite scalar spectral identity integrates exactly, and the
finite Weyl gap satisfies
\begin{align}
    \sum_gD(A_g\Vert B_g)-D(A\Vert B)
        &=
        \int_0^\infty
        \frac{\operatorname{Def}_A(t)+t\operatorname{Def}_B(t)}
             {1+t}\,dt.
    \label{eq:cpsv16_ssa_equality_hayashi_markov_positive_definite_gap_integral}
\end{align}
The normalization is essential. Saturation of the finite Weyl Jensen
inequality gives $d_R^{-2}$ times each relative entropy, whereas
Eq.~(\ref{eq:cpsv16_ssa_equality_hayashi_markov_positive_definite_gap_integral})
uses the weighted matrices $A_g$ and $B_g$. Positive homogeneity,
\begin{align}
    D(cA\Vert cB)
        &=cD(A\Vert B),
        \qquad c>0,
    \label{eq:cpsv16_ssa_equality_hayashi_markov_relative_entropy_homogeneity}
\end{align}
identifies the two expressions. Equality under the partial trace therefore
makes the weighted gap vanish without an additional gap hypothesis.

Both defects are nonnegative and continuous on $(0,\infty)$, and the full
integrand in
Eq.~(\ref{eq:cpsv16_ssa_equality_hayashi_markov_positive_definite_gap_integral})
is integrable. A zero gap gives almost-everywhere vanishing, and continuity
then gives pointwise vanishing for every $t>0$. Since $t/(1+t)>0$, the
source-$B$ defect vanishes, and
Eq.~(\ref{eq:cpsv16_ssa_equality_hayashi_markov_common_resolvent}) follows.%
\footnote{Formal declarations:
\leanid{Matrix.quantumRelativeEntropy_smul_posDef},
\leanid{Matrix.weylRelativeEntropyGap_eq_zero_of_partialTraceRight_eq},
\leanid{Matrix.quantumRelativeEntropy_resolvent_integral},
\leanid{Matrix.weyl_relativeEntropy_gap_integral},
\leanid{Matrix.weyl_sourceB_defect_eq_zero_of_gap_eq_zero},
\leanid{Matrix.weyl_resolvent_eq_of_relativeEntropy_gap_eq_zero}, and
\leanid{Matrix.weyl_resolvent_eq_of_partialTraceRight_eq} in
\path{TNLean/Analysis/RelativeEntropyResolventIntegral.lean} and
\path{TNLean/Channel/Schwarz/WeylRelativeEntropyIntegral.lean}.}

Define the relative modular matrix
\begin{align}
    \Delta_{A,B}
        &=A\otimes(B^{-1})^{\mathsf T}.
    \label{eq:cpsv16_ssa_equality_hayashi_markov_relative_modular_matrix}
\end{align}
The source-$B$ resolvent factors through this matrix. Equality of the shifted
resolvents on the vectorized identity passes through the positive square-root
functional calculus and gives
\begin{align}
    \sqrt{\rho} \sigma^{-1/2}
        &=\sqrt{\overline\rho} \overline\sigma^{-1/2}.
    \label{eq:cpsv16_ssa_equality_hayashi_markov_positive_sqrt_ratio}
\end{align}
Taking adjoint products and cancelling $\sqrt{\sigma}$ gives
\begin{align}
    \sqrt{\sigma}
    \overline\sigma^{-1/2}
    \overline\rho
    \overline\sigma^{-1/2}
    \sqrt{\sigma}
        &=\rho.
    \label{eq:cpsv16_ssa_equality_hayashi_markov_positive_petz_sandwich}
\end{align}
This is the positive-square-root specialization of the relative-modular
analytic-continuation argument in the Jen\v{c}ov\'a--Ruskai preprint
arXiv:0903.2895v4, lines 658--680. The formal proof uses the equivalent
L\"owner integral representation for the power $1/2$.\footnote{Formal
declarations:
\leanid{Matrix.sqrt_mulVec_eq_of_resolvent_mulVec_eq},
\leanid{Matrix.mulVec_sqrt_mulVec_eq_of_mulVec_resolvent_mulVec_eq},
\leanid{Matrix.sourceB_resolvent_eq_relativeModular}, and
\leanid{Matrix.relativeModular_sqrt_mulVec_vec_one} in
\path{TNLean/Analysis/ResolventFunctionalCalculus.lean}; and
\leanid{Matrix.weyl_sqrt_ratio_eq_of_partialTraceRight_eq},
\leanid{Matrix.weyl_identity_sandwich_of_partialTraceRight_eq}, and
\leanid{Matrix.partialTraceRightPetzMap_eq_of_relativeEntropy_eq_posDef} in
\path{TNLean/Channel/Schwarz/WeylRelativeEntropyIntegral.lean}.}
Thus partial-trace equality implies raw Petz recovery for positive definite
inputs. The support-domain argument below extends this result to singular
inputs.

\section{Singular-support resolvent argument}

The Jen\v{c}ov\'a--Ruskai preprint arXiv:0903.2895v4, lines 717--720, extends
the definition and convexity argument to positive semidefinite pairs $(A,B)$
satisfying $\ker B\subseteq\ker A$. Its singular equality argument, at
lines 766--793, obtains the common relative-modular resolvent on the reference
support and then applies functional calculus. These relative-entropy equality
methods are reviewed in Ref.~\cite{Cirac2021Matrix}.

The support-domain spectral integral is now proved. For positive semidefinite
$A,B$ with $\ker B\subseteq\ker A$, the finite spectral integrand is
integrable on $(0,\infty)$ and satisfies
\begin{align}
    \int_0^\infty
        \operatorname{Spec}_{A,B}(t)\,dt
        &=D(A\Vert B).
    \label{eq:cpsv16_ssa_equality_hayashi_markov_support_spectral_integral}
\end{align}
This proves the spectral content of the formulas denoted
$(\mathrm{J1})$, $(\mathrm{intspec})$, and $(\mathrm{intAB})$ in that
preprint, together with its support convention at lines 717--720.%
\footnote{Formal declarations:
\leanid{Matrix.quantumRelativeEntropy_spectral} and
\leanid{Matrix.supportRelativeEntropySpectralIntegrand_integrableOn_and_integral_eq_quantumRelativeEntropy}
in \path{TNLean/Analysis/RelativeEntropySupportIntegral.lean}, and
\leanid{Matrix.supportRelativeEntropyLeftRightIntegrand_eq_spectral} and
\leanid{Matrix.rawSupportSourceAQuadratic_sub_trace_add_sourceB_eq_spectral}
in \path{TNLean/Analysis/RelativeEntropySupportLeftRightQuadratic.lean}.}

The spectral integrand agrees pointwise with the support-projected left--right
quadratic form. Its first source term is $AP_B$, where $P_B$ is the support
projection of $B$. The kernel inclusion implies
\begin{align}
    AP_B
        &=A.
    \label{eq:cpsv16_ssa_equality_hayashi_markov_source_support_absorption}
\end{align}
Hence this term equals the quadratic form with unprojected source $A$.

Define the generalized inverse
\begin{align}
    B^+
        &=\left(B^{-1/2}_{\mathrm{supp}}\right)^2.
    \label{eq:cpsv16_ssa_equality_hayashi_markov_reference_generalized_inverse}
\end{align}
For a single pair, the support-generalized-inverse source-$B$ solution equals
the projected shifted relative-modular vector
\begin{align}
    &(\Id\otimes P_B^{\mathsf T})
    \left(t\Id+A\otimes(B^+)^{\mathsf T}\right)^{-1}
    \operatorname{vec}(\Id^{\mathsf T}).
    \label{eq:cpsv16_ssa_equality_hayashi_markov_projected_modular_vector}
\end{align}
The corresponding source-$B$ quadratic forms are therefore equal.%
\footnote{Formal declarations:
\leanid{Matrix.supportLeftRightSupportInv_mulVec_sourceB_eq_projected_relativeModular}
and
\leanid{Matrix.supportSourceBQuadratic_eq_projected_relativeModular} in
\path{TNLean/Channel/Schwarz/SupportLeftRightRelativeModular.lean}.}

For a finite family $(A_i,B_i)$ satisfying
$\ker B_i\subseteq\ker A_i$, both fixed-parameter defects are nonnegative:
\begin{align}
    \sum_iQ^A_{A_i,B_i}(t)
        -Q^A_{\sum_iA_i,\sum_iB_i}(t)
        &\geq 0,
    \label{eq:cpsv16_ssa_equality_hayashi_markov_source_a_defect_nonnegative}\\
    \sum_iQ^B_{A_i,B_i}(t)
        -Q^B_{\sum_iA_i,\sum_iB_i}(t)
        &\geq 0.
    \label{eq:cpsv16_ssa_equality_hayashi_markov_source_b_defect_nonnegative}
\end{align}
The source in
Eq.~(\ref{eq:cpsv16_ssa_equality_hayashi_markov_source_a_defect_nonnegative})
is the unprojected matrix $A_i$. The source-$B$ inequality in
Eq.~(\ref{eq:cpsv16_ssa_equality_hayashi_markov_source_b_defect_nonnegative})
requires no kernel inclusion. The local kernel inclusions imply the
corresponding inclusion for the summed pair.\footnote{Formal declarations:
\leanid{Matrix.rawSupportSourceAQuadratic_sum_sub_nonneg},
\leanid{Matrix.supportSourceBQuadratic_sum_sub_nonneg},
\leanid{Matrix.supportLeftRightSupportInv_mulVec_sourceB_eq_of_defect_eq_zero},
and
\leanid{Matrix.supportRightProj_mul_sumSupportRightProj_eq} in
\path{TNLean/Channel/Schwarz/SupportSourceDefect.lean}.}

The finite-family relative-entropy gap is the integral of these two defects:
\begin{align}
    \sum_iD(A_i\Vert B_i)
        -D\!\left(\sum_iA_i\middle\Vert\sum_iB_i\right)
        &=
        \int_0^\infty
        \frac{\operatorname{Def}_A(t)+t\operatorname{Def}_B(t)}
             {1+t}\,dt.
    \label{eq:cpsv16_ssa_equality_hayashi_markov_support_family_gap}
\end{align}
The gap integrand is continuous and integrable on $(0,\infty)$.%
\footnote{Formal declarations:
\leanid{Matrix.supportRelativeEntropyFamilyGapIntegrand},
\leanid{Matrix.supportRelativeEntropyFamilyGapIntegrand_integrableOn_and_integral_eq},
and
\leanid{Matrix.supportRelativeEntropyFamilyGapIntegrand_continuousOn} in
\path{TNLean/Channel/Schwarz/SupportRelativeEntropyGap.lean}.}
Equations~(\ref{eq:cpsv16_ssa_equality_hayashi_markov_source_a_defect_nonnegative})
and
(\ref{eq:cpsv16_ssa_equality_hayashi_markov_source_b_defect_nonnegative}) make
this integrand nonnegative. If the relative entropy of the summed pair equals
the sum of the individual relative entropies, the integral vanishes.
Almost-everywhere vanishing and continuity then imply pointwise vanishing for
every $t>0$, including vanishing of the source-$B$ defect.%
\footnote{Formal declarations:
\leanid{Matrix.supportRelativeEntropyFamilyGapIntegrand_nonneg} and
\leanid{Matrix.supportSourceBDefect_eq_zero_of_relativeEntropy_sum_eq} in
\path{TNLean/Channel/Schwarz/SupportRelativeEntropyGap.lean}.}

The residual identity converts the real-valued defect equality into a common
support-generalized-inverse solution. The local right-support projection
absorbs both the support projection of the local left--right operator and the
right-support projection of the summed reference matrix. The local and summed
shifted relative-modular resolvents therefore agree after the local
right-support projection for every summand and every $t>0$.%
\footnote{Formal declarations:
\leanid{Matrix.supportRightProj_mul_supportLeftRightSupportProj_eq} in
\path{TNLean/Channel/Schwarz/SupportLeftRightRelativeModular.lean}, and
\leanid{Matrix.supportRelativeModular_resolvent_mulVec_eq_of_relativeEntropy_sum_eq}
in \path{TNLean/Channel/Schwarz/SupportRelativeEntropyEquality.lean}.}

For the Weyl family, now use the unweighted conjugates
\begin{align}
    A_g
        &=U_g\rho U_g^\dagger,
    \label{eq:cpsv16_ssa_equality_hayashi_markov_unweighted_weyl_state}\\
    B_g
        &=U_g\sigma U_g^\dagger.
    \label{eq:cpsv16_ssa_equality_hayashi_markov_unweighted_weyl_reference}
\end{align}
Positive-scalar homogeneity on the support domain converts weighted Jensen
equality into
\begin{align}
    D\!\left(\sum_gA_g\middle\Vert\sum_gB_g\right)
        &=\sum_gD(A_g\Vert B_g).
    \label{eq:cpsv16_ssa_equality_hayashi_markov_unweighted_weyl_gap_zero}
\end{align}
Unitary conjugation transports the local kernel inclusions, and positivity
transports them to the summed pair. The generic finite-family theorem therefore
gives the common projected relative-modular resolvent for every Weyl summand.
At the zero Weyl index, the local pair is $(\rho,\sigma)$.\footnote{Formal
declarations:
\leanid{Matrix.quantumRelativeEntropy_smul_support} in
\path{TNLean/Analysis/RelativeEntropySupportIntegral.lean}, and
\leanid{Matrix.weyl_supportRelativeModular_resolvent_mulVec_eq_of_partialTraceRight_eq}
in
\path{TNLean/Channel/Schwarz/WeylSupportRelativeEntropyEquality.lean}.}
This result supplies the analytic input for the support square-root-ratio step.
It does not yet assert the square-root ratio, the Petz sandwich, or recovery.

\section{Support square-root ratio and recovery}

Consider positive semidefinite pairs $(A,B)$ and $(C,D)$. Let $P_B$ be the
support projection of $B$, and define
\begin{align}
    B^+
        &=\left(B^{-1/2}_{\mathrm{supp}}\right)^2,
    \label{eq:cpsv16_ssa_equality_hayashi_markov_b_generalized_inverse}\\
    D^+
        &=\left(D^{-1/2}_{\mathrm{supp}}\right)^2.
    \label{eq:cpsv16_ssa_equality_hayashi_markov_d_generalized_inverse}
\end{align}
Suppose that, for every $t>0$, the projected shifted resolvents satisfy
\begin{align}
    &(\Id\otimes P_B^{\mathsf T})
    \left(t\Id+A\otimes(B^+)^{\mathsf T}\right)^{-1}
    \operatorname{vec}(\Id^{\mathsf T})
        \notag\\
    &\qquad=
    (\Id\otimes P_B^{\mathsf T})
    \left(t\Id+C\otimes(D^+)^{\mathsf T}\right)^{-1}
    \operatorname{vec}(\Id^{\mathsf T}).
    \label{eq:cpsv16_ssa_equality_hayashi_markov_projected_resolvent_equality}
\end{align}
Then the square-root ratios agree on the same support:
\begin{align}
    \left(\sqrt A\,B^{-1/2}_{\mathrm{supp}}\right)P_B
        &=
        \left(\sqrt C\,D^{-1/2}_{\mathrm{supp}}\right)P_B.
    \label{eq:cpsv16_ssa_equality_hayashi_markov_support_sqrt_ratio}
\end{align}
The vectorized restriction has the orientation
\begin{align}
    (\Id\otimes P_B^{\mathsf T})
    \operatorname{vec}(M^{\mathsf T})
        &=
        \operatorname{vec}\!\left((MP_B)^{\mathsf T}\right).
    \label{eq:cpsv16_ssa_equality_hayashi_markov_vectorized_support_orientation}
\end{align}
\footnote{Formal declarations:
\leanid{Matrix.one_kronecker_transpose_mulVec_vec_transpose},
\leanid{Matrix.supportRelativeModular_sourceB_solution},
\leanid{Matrix.supportRelativeModular_sqrt_mulVec_vec_one}, and
\leanid{Matrix.supportRelativeModular_sqrt_ratio_eq_of_resolvent_mulVec_eq} in
\path{TNLean/Channel/Schwarz/SupportRelativeModular.lean}.}

The projection in
Eq.~(\ref{eq:cpsv16_ssa_equality_hayashi_markov_projected_resolvent_equality})
is essential. Equality on the ambient vectorized identity would force equality
of the two reference support projections because each resolvent contributes
$t^{-1}$ on the complement of its reference support. This conclusion is false
in the intended data-processing application. For example, if $\rho=\sigma$ is
a rank-one maximally entangled state, partial-trace equality is automatic, but
the comparison pair formed from its marginal is full rank. The singular
equality theorem therefore does not imply ambient resolvent equality.

The canonical source-side projected resolvent solves its singular left--right
equation. After applying $\Id\otimes P_B^{\mathsf T}$, multiplication by the
operator
\begin{align}
    A\otimes\Id+t(\Id\otimes B^{\mathsf T})
    \label{eq:cpsv16_ssa_equality_hayashi_markov_singular_left_right_operator}
\end{align}
gives $\operatorname{vec}(B^{\mathsf T})$. This uses only the
generalized-inverse identity
\begin{align}
    \left(B^{-1/2}_{\mathrm{supp}}\right)^2B
        &=P_B.
    \label{eq:cpsv16_ssa_equality_hayashi_markov_generalized_inverse_identity}
\end{align}
The one-pair identity also identifies this vector with the
support-generalized-inverse solution of the left--right equation and therefore
identifies the two source-$B$ quadratic forms.

Together with the integral and continuity argument, equality of relative
entropies makes the real-valued source-$B$ defect vanish at every positive
parameter. Nonnegativity of the complex quadratic residuals makes their
imaginary parts vanish, so the real equality yields the finite-family common
projected resolvent theorem. The kernel inclusion required for the unprojected
source-$A$ inequality remains an exact hypothesis of this theorem. The finite
Weyl specialization supplies the projected resolvent hypothesis for the
support functional calculus. The zero Weyl summand, the unweighted-sum
normalization, and
Eq.~(\ref{eq:cpsv16_ssa_equality_hayashi_markov_support_sqrt_ratio}) then give
the singular inverse-square-root sandwich and raw Petz recovery.%
\footnote{Formal declarations:
\leanid{Matrix.unweightedWeylSum_eq_smul_partialTraceRight_kronecker} in
\path{TNLean/Channel/Schwarz/WeylSupportRelativeEntropyEquality.lean};
\leanid{Matrix.weyl_support_sqrt_ratio_eq_of_partialTraceRight_eq},
\leanid{Matrix.weyl_support_identity_sandwich_of_partialTraceRight_eq}, and
\leanid{Matrix.partialTraceRightPetzMap_eq_of_relativeEntropy_eq} in
\path{TNLean/Channel/Schwarz/WeylSupportPetzRecovery.lean}.}

Affine regularization alone cannot prove the equality case. Equality at the
limiting pair implies only that the nonnegative regularized gaps converge to
zero; it does not imply that a gap at any positive regularization parameter is
zero. The positive-definite equality theorem therefore cannot be applied
pointwise to the regularized pairs without a separate equality-preservation
theorem.

For the identity Weyl summand, write
\begin{align}
    \overline X
        &=(\tr_R X)\otimes d_R^{-1}\Id_R.
    \label{eq:cpsv16_ssa_equality_hayashi_markov_identity_weyl_average}
\end{align}
For every positive semidefinite $\tau$ on the left factor,
\begin{align}
    \left(\tau\otimes d_R^{-1}\Id_R\right)^{-1/2}_{\mathrm{supp}}
        &=
        \sqrt{d_R}
        \left(\tau^{-1/2}_{\mathrm{supp}}\otimes\Id_R\right),
    \label{eq:cpsv16_ssa_equality_hayashi_markov_mixed_support_inverse_sqrt}\\
    \overline\sigma^{-1/2}_{\mathrm{supp}}
    \overline\rho
    \overline\sigma^{-1/2}_{\mathrm{supp}}
        &=
        \left((\tr_R\sigma)^{-1/2}_{\mathrm{supp}}
            (\tr_R\rho)
            (\tr_R\sigma)^{-1/2}_{\mathrm{supp}}\right)\otimes\Id_R.
    \label{eq:cpsv16_ssa_equality_hayashi_markov_weyl_sandwich_reduction}
\end{align}
The two factors $\sqrt{d_R}$ in the second line cancel the coefficient
$d_R^{-1}$ in $\overline\rho$. Consequently, the operator identity
\begin{align}
    \sqrt{\sigma}
    \overline\sigma^{-1/2}_{\mathrm{supp}}
    \overline\rho
    \overline\sigma^{-1/2}_{\mathrm{supp}}
    \sqrt{\sigma}
        &=\rho
    \label{eq:cpsv16_ssa_equality_hayashi_markov_support_petz_sandwich}
\end{align}
implies
\begin{align}
    \mathcal R_\sigma(\tr_R\rho)
        &=\rho
    \label{eq:cpsv16_ssa_equality_hayashi_markov_raw_petz_recovery}
\end{align}
through the support Petz formula
(\ref{eq:cpsv16_ssa_equality_hayashi_markov_petz_support_map}), corresponding
to Theorem~3, Equation~(8), of the Hayden--Jozsa--Petz--Winter argument reviewed
in Ref.~\cite{Cirac2021Matrix}.\footnote{Formal declarations:
\leanid{Matrix.PosSemidef.supportInvSqrt_kronecker_maximallyMixed},
\leanid{Matrix.PosSemidef.supportInvSqrt_kronecker_maximallyMixed_mul_mul}, and
\leanid{Matrix.partialTraceRightPetzMap_eq_of_weyl_identity_sandwich} in
\path{TNLean/Channel/Schwarz/PetzEqualityPrerequisites.lean}.}
The implication is conditional on
Eq.~(\ref{eq:cpsv16_ssa_equality_hayashi_markov_support_petz_sandwich}). The
finite Weyl argument now proves that identity from partial-trace equality for
positive semidefinite inputs satisfying
$\ker\sigma\subseteq\ker\rho$. No corresponding claim is made for an arbitrary
convex ensemble.

The singular-support equality theorem in Weyl coordinates is therefore
\begin{align}
    D(\rho\Vert\sigma)
        &=D(\tr_R\rho\Vert\tr_R\sigma)
    \label{eq:cpsv16_ssa_equality_hayashi_markov_partial_trace_dpi_equality}\\
    &\Longrightarrow
    \mathcal R_\sigma(\tr_R\rho)=\rho.
    \label{eq:cpsv16_ssa_equality_hayashi_markov_partial_trace_raw_recovery}
\end{align}
This is the raw support-map identity used in the forward implication of
Theorem~3 in the Hayden--Jozsa--Petz--Winter argument reviewed in
Ref.~\cite{Cirac2021Matrix}. The proof does not infer equality for positive
affine regularizations. Instead, the support-restricted finite-family equality
theorem gives the common projected Weyl resolvent, support functional calculus
gives the projected square-root ratio, and the kernel inclusion reduces its
Gram identity to
Eq.~(\ref{eq:cpsv16_ssa_equality_hayashi_markov_support_petz_sandwich}).

Product-coordinate covariance transports this identity to arbitrary finite
tensor factors. For density operators, the support condition and
Eq.~(\ref{eq:cpsv16_ssa_equality_hayashi_markov_completed_raw_agreement}) show
that the completed channel agrees with the raw map on $\tr_R\rho$. Thus,
\begin{align}
    D(\rho\Vert\sigma)
        &=D(\tr_R\rho\Vert\tr_R\sigma)
    \label{eq:cpsv16_ssa_equality_hayashi_markov_general_dpi_equality}\\
    &\Longrightarrow
    \widehat{\mathcal R}_\sigma(\tr_R\rho)=\rho.
    \label{eq:cpsv16_ssa_equality_hayashi_markov_completed_recovery}
\end{align}
The complementary preparation term in $\widehat{\mathcal R}_\sigma$ is the
chosen trace-preserving extension; it is not part of Equation~(8) in the
Hayden--Jozsa--Petz--Winter argument reviewed in
Ref.~\cite{Cirac2021Matrix}. The corresponding formal declarations are
\leanid{Matrix.partialTraceRightPetzMap_submatrix_prod_equiv},
\leanid{Matrix.partialTraceRightPetzMap_eq_of_relativeEntropy_eq_general_support},
and
\leanid{Matrix.partialTraceRightPetzChannel_recovery_of_dpi_equality}.

\section{Application to strong subadditivity}

For the strong-subadditivity pair, set
\begin{align}
    \rho
        &=\rho_{ABC},
    \label{eq:cpsv16_ssa_equality_hayashi_markov_ssa_state}\\
    \sigma
        &=\rho_A\otimes\rho_{BC}.
    \label{eq:cpsv16_ssa_equality_hayashi_markov_ssa_reference}
\end{align}
The required product-marginal support inclusion is automatic. The formal
development therefore proves the recovery statement corresponding to
Equation~(11) of the Hayden--Jozsa--Petz--Winter argument reviewed in
Ref.~\cite{Cirac2021Matrix}:
\begin{align}
    \rho_{ABC}
        &=
        \left(\operatorname{id}_A\otimes
            \widehat{\mathcal R}_{\rho_{BC}}\right)(\rho_{AB}).
    \label{eq:cpsv16_ssa_equality_hayashi_markov_hjpw_recovery_identity}
\end{align}
The relevant formal declarations are
\leanid{Matrix.product_marginal_support},
\leanid{Matrix.partialTraceRightPetzMap_product_marginal_recovery_of_isSSAEquality},
\leanid{Matrix.partialTraceRightPetzChannel_traceA_ABC_isKrausCPTP}, and
\leanid{Matrix.idTensor_partialTraceRightPetzChannel_traceC_ABC_eq_of_isSSAEquality}.
The completed local Petz map is completely positive and trace preserving. If
$\rho_A$ is singular, this is a supported-input identity rather than a global
factorization of TNLean's generic completed product-reference channel.

The full-support invariant-state decomposition, the preserving-operation block
action used in Theorem~6 of the Hayden--Jozsa--Petz--Winter argument, the
joint-support reduction, and its application to the invariant conditional
family are formalized. The relevant declarations are
\leanid{Matrix.exists_markovBipartiteBlockForm_jointSupport},
\leanid{Matrix.exists_ambientMarkovBipartiteBlockForm},
\leanid{Matrix.exists_markovDilationBlockForm},
\leanid{Matrix.MarkovDilationBlockForm.ambient_tripartite_block_form}, and
\leanid{Matrix.exists_markovTripartiteBlockForm}. These results produce the
positive unnormalized factors $\omega_j$, prove that their traces sum to one,
produce the trace-one supported sector states $\widehat\rho_j$, and establish
Eq.~(\ref{eq:cpsv16_ssa_equality_hayashi_markov_ambient_hjpw}) with zero
complementary blocks.

The forward characterization theorem
\leanid{Matrix.hayashi_ssa_equality_characterization_forward} then sets
$p_j=\operatorname{Re}\tr(\omega_j)$, reverses the fibre order, uses the
adjoint middle-system unitary, and applies total normalization. Every supported
right factor remains the corresponding sector output state, including in a
zero-weight supported sector. Such a sector may use a filler only on the left;
a complementary sector may use fillers on both sides.

The singular-support recovery scope is unchanged. Equality in data processing
gives raw support-Petz recovery, and the chosen completed local channel agrees
with the raw map on the Petz output marginal. No tensor-product factorization
is claimed for the generic completed singular product-reference channel away
from the supported input.

\section{Verdict}

\emph{Closed.} Both directions of the strong-subadditivity equality
characterization are proved, and the resulting biconditional is axiom-free.
The forward proof derives the ambient Hayden--Jozsa--Petz--Winter blocks,
chooses $p_j=\operatorname{Re}\tr(\omega_j)$ with nonnegative weights summing
to one, retains every trace-one sector output state on supported sectors, and
introduces normalized fillers only in sectors whose final weight is zero. The
middle fibres have the order $B_j^L\otimes B_j^R$, and the Hayashi unitary has
the adjoint orientation required by
Eq.~(\ref{eq:cpsv16_ssa_equality_hayashi_markov_ambient_hjpw}).

This closure does not strengthen the completed Petz map away from the support.
The proved source-facing result remains raw support-map recovery together with
agreement of the selected completed local channel on the supported marginal.

\begin{thebibliography}{9}

\bibitem{Ruskai2002}
  M.~B.~Ruskai,
  \emph{Inequalities for quantum entropy: a review with conditions for
  equality},
  Journal of Mathematical Physics \textbf{43} (2002), 4358--4375.

\bibitem{HaydenJozsaPetzWinter2004}
  P.~Hayden, R.~Jozsa, D.~Petz, and A.~Winter,
  \emph{Structure of states which satisfy strong subadditivity of quantum
  entropy with equality},
  Communications in Mathematical Physics \textbf{246} (2004), 359--374.

\bibitem{Hayashi2006}
  M.~Hayashi,
  \emph{Quantum Information: An Introduction},
  Springer, 2006; Theorem 5.24.

\end{thebibliography}

\end{document}
